\chapter{Case Studies I}\label{ch:fm:case-studies-i}

On 23 September 1998 fourteen banks and securities firms agreed to recapitalise a hedge fund, Long-Term Capital Management, and five days later
they put in about \$3.6 billion, 90\% of what the fund was then worth. The fund had begun the year with \$4.67 billion of capital and a balance
sheet 28 times as large; by the end of August it had lost almost half, and by the recapitalisation almost all. Every number in that paragraph was
the result of decisions the firm took long before: how much leverage to run, on what funding terms, in which trades, and alongside whom. This chapter
and the next read three such failures from the public record, and rerun each with one decision changed.

\section{How to read a failure}

A failure seen from outside is a story, and stories favour villains and bad luck. Seen from inside a firm, it is a sequence of decisions, each of
which looked reasonable with the information of its time: a limit set, a funding term accepted, a position allowed to grow, a choice to sell or
hold. The useful reading keeps to what the public record establishes, reconstructs the numbers from it, and asks which decision, changed, would
have changed the outcome by how much. It does not re-litigate what regulators and courts found, and it does not attribute motives the record does
not state.

Two kinds of liquidity run through all three cases, and Book 1's liquidity spiral is the mechanism that joins them.

\begin{definition}[Funding liquidity risk, market liquidity risk]\label{def:fm:case-studies-i:liquidity}
\emph{Funding liquidity risk}\index{funding liquidity risk} is the risk that a firm cannot meet its payments and margin calls when they fall due,
because its lenders raise haircuts, shorten terms or withdraw, or because its own losses consume its cash. \emph{Market liquidity
risk}\index{market liquidity risk} is the risk that a firm cannot sell or hedge a position quickly without moving its price against itself, because
the position is large against the market's volume or because other holders are selling at the same time.
\end{definition}

A firm with ample funding can wait for a market to recover; a firm whose positions can be sold at once needs little funding patience. The failures
that end firms combine the two: losses trigger margin calls, the calls force sales, the sales move prices, and the new prices trigger more calls
(\cref{fig:fm:case-studies-i:spiral}).

\begin{omfigure}
\begin{tikzpicture}[font=\small,box/.style={draw=omDef,fill=omDef!8,rounded corners,align=center,minimum width=3.3cm,minimum height=0.95cm},
  arr/.style={-{Stealth[length=2.2mm]},thick,omDef}]
\node[box] (loss) at (0,2.2) {losses on\\positions};
\node[box] (call) at (6.4,2.2) {margin calls,\\higher haircuts};
\node[box] (sell) at (6.4,0) {forced sales};
\node[box] (price) at (0,0) {prices move\\against holders};
\draw[arr] (loss) -- node[above,font=\scriptsize] {funding liquidity} (call);
\draw[arr] (call) -- (sell);
\draw[arr] (sell) -- node[below,font=\scriptsize] {market liquidity} (price);
\draw[arr] (price) -- (loss);
\node[draw=omThm,dashed,rounded corners,align=center,font=\scriptsize,text width=3.1cm] at (10.5,1.1) {others holding the\\same positions mark\\them at the new prices};
\draw[-{Stealth[length=2mm]},omThm,dashed] (8.85,1.1) -- (8.15,0.35);
\end{tikzpicture}
\omcaption{The loop that joins the two kinds of liquidity risk: losses raise calls, calls force sales, sales move prices, and moved prices are new
losses, for the seller and for everyone who holds the same positions.}\label{fig:fm:case-studies-i:spiral}
\end{omfigure}

\begin{definition}[Days to liquidate]\label{def:fm:case-studies-i:days}
The \emph{days to liquidate}\index{days to liquidate} a position is its size divided by the volume the firm can trade each day without an
unacceptable market impact, usually a stated share of the average daily volume:
\[ D=\frac{Q}{\pi\,V}, \]
for a position of $Q$ units, an average daily volume of $V$ units and a participation rate $\pi$ (Book 10's participation).
\end{definition}

\omterm{def:fm:case-studies-i:days}{Days to liquidate} is the bridge between a position's size and a risk model's horizon. A one-day value at risk (Book 6) assumes the position can be
exited within the day at prices near the model's; a position of twenty \omterm{def:fm:case-studies-i:days}{days to liquidate} carries twenty days of price risk, and its own sales move the
price on each of them.

\section{1998: leverage and the terms of funding}

The public record here is the report of the General Accounting Office (now the Government Accountability Office) of October 1999. At the end of 1997 the fund
returned about \$2.7 billion of capital to its investors, leaving net assets of \$4.67 billion; its balance-sheet leverage, 17 to 1 at the end of 1994,
stood at 28 to 1. On 31 August 1998 it also held derivatives of about \$1.4 trillion notional off its balance sheet. From January to September 1998 it
lost almost 90\% of its capital: 44\% in August alone, after the Russian debt moratorium of mid-August sent investors towards quality and liquidity, and most of what remained in September.

\begin{dated}{2026-09}{The public record of the chapter's cases}\label{dat:fm:case-studies-i:record}
\textbf{Long-Term Capital Management} (GAO/GGD-00-3, October 1999): net assets \$4.67 billion at the end of 1997, after about \$2.7 billion returned
to investors; balance-sheet leverage 28 to 1 at the end of 1997; \$2.3 billion of net assets on 31 August 1998, after a 44\% loss in August; on 23
September 1998, 14 banks and securities firms agreed to recapitalise it, and on 28 September contributed about \$3.6 billion, 90\% of its net asset
value. \textbf{Amaranth Advisors} (US Senate Permanent Subcommittee on Investigations, hearing record S.~Hrg.~110-235, 2007): at times in the summer of
2006 it held about 40\% of all outstanding NYMEX natural gas contracts for the winter of 2006--2007, 75\% of November 2006 and 60\% of January and
March 2007, and at times 100\,000 contracts in a month; from the end of August to mid-September 2006 its natural gas positions lost more than \$2
billion, and it liquidated its natural gas portfolio. \textbf{Global Equity Opportunities} (Goldman Sachs, Form 8-K, 13 August 2007): a quantitative
long/short equity fund with a net asset value of about \$3.6 billion received a \$3 billion equity investment from Goldman Sachs and other investors.
\end{dated}

A balance sheet 28 times capital means that a fall of $1/28=3.6\%$ in the value of the assets, net of hedges, wipes out the capital. The fund's
positions were mostly spread trades, long one bond or swap and short a close relative, whose net value moves far less than either leg; the leverage
was sized to that small net risk. What the ratio measures is how little room the funding leaves when the spreads move together and far.

The GAO report states why the Federal Reserve facilitated a private recapitalisation: it judged that a rapid liquidation of the fund's
positions, and of related positions held by other market participants, might pose a significant threat to markets that were already unsettled. The
creditors held the same trades, or were the counterparties to them; a forced sale by the fund would have marked their books too.

\section{Tutorial: three reconstructions}\label{tut:fm:case-studies-i:tut}

\begin{tutorial}
\textbf{Goal.} Reconstruct from the public figures the 1998 fund's capital through the year and rerun it at half the leverage; reconstruct the
natural-gas fund's share of open interest and its \omterm{def:fm:case-studies-i:days}{days to liquidate}; rerun Book 7's crowded unwind for a fund that sells and one that holds.
\textbf{End state:} the three charts \cref{fig:fm:case-studies-i:ltcm,fig:fm:case-studies-i:amaranth,fig:fm:case-studies-i:august}.
\begin{tutsteps}
\item \textbf{Capital by period.} From \$2.3 billion at the end of August after a 44\% loss, capital at the end of July was $2.3/0.56=\$4.11$
billion. The \$3.6 billion was 90\% of the net asset value on 28 September, so the existing investors' share was $3.6/0.9-3.6=\$0.40$ billion.
\item \textbf{Returns on capital.} January to July $4.11/4.67-1=-12.1\%$; August $-44\%$; September to the 28th $0.40/2.3-1=-82.6\%$.
\item \textbf{Half the leverage.} If the assets' returns are unchanged, halving the leverage halves each period's return on capital:
$4.67\times(1-0.061)(1-0.22)(1-0.413)=\$2.01$ billion on 28 September (\texttt{firm.casebook.capital\_path}).
\item \textbf{Concentration.} The subcommittee's shares of open interest, and \omterm{def:fm:case-studies-i:days}{days to liquidate} for 100\,000 contracts at 20\% of an illustrative
average daily volume of 25\,000 contracts in the month's contract.
\item \textbf{The crowded exit.} Three funds with long-short books of \$30 billion gross at six times leverage and an overlap of 0.94; one sells half
its book over five days, the others hold (\texttt{firm.casebook.quant\_unwind}, on \texttt{firm.capacity.unwind}).
\end{tutsteps}
\end{tutorial}

\omcode{code/firm/casebook/firm_casebook.py}{35}{56}{Returns on capital derived from the public capital figures, the capital path at a scaled leverage,
and days to liquidate.}\label{lst:fm:case-studies-i:recon}

\begin{omfigure}
\begin{tikzpicture}
\begin{axis}[width=10.5cm,height=5.4cm,ybar,bar width=17pt,ymin=0,ymax=5.4,ylabel={\small capital (\$ billion)},xtick={0,1,2,3},
  xticklabels={end 1997,end July,end August,28 September},xmin=-0.6,xmax=3.6,tick label style={font=\footnotesize},grid=major,grid style={gray!20},
  table/col sep=comma,nodes near coords,nodes near coords style={font=\tiny,/pgf/number format/fixed,/pgf/number format/precision=2},
  legend style={font=\scriptsize,at={(0.5,-0.22)},anchor=north,/tikz/every even column/.append style={column sep=8pt}},legend columns=2,
  legend cell align=left]
\addplot[fill=omDef!60,draw=omDef] table[x=pos,y=actual]{figdata/desk/28-case-studies-i/ltcm.csv};
\addplot[fill=omThm!45,draw=omThm] table[x=pos,y=half]{figdata/desk/28-case-studies-i/ltcm.csv};
\legend{reconstructed from the record,at half the leverage}
\end{axis}
\end{tikzpicture}
\omcaption{The 1998 fund's capital, reconstructed from the GAO figures, against the same returns on its assets at half the leverage: \$0.40 billion
of the investors' capital left on 28 September against \$2.01 billion. Data: \texttt{fm\_cases1.ltcm}.}\label{fig:fm:case-studies-i:ltcm}
\end{omfigure}

The reconstruction leaves the investors with \$0.40 billion of their \$4.67 billion, a loss of 91.4\%, before the recapitalisation; the net asset
value of \$3.81 billion reported for 30 September includes the new money. At half the leverage the same asset returns would have left \$2.01 billion,
43\% of the capital, and no need for rescue. That counterfactual is generous in one way and harsh in another. It assumes the assets would have
fallen as far, when part of September's fall was the market trading against a fund known to be in distress, which a smaller fund would not have been;
and it gives the smaller fund no credit for fewer counterparties calling margin at once. It also takes nothing from the returns of the years before,
which half the leverage would have halved.

\begin{remark}[Leverage and the term of funding]\label{rem:fm:case-studies-i:term}
Leverage alone does not end a fund; losses end it when its lenders can demand margin faster than its positions recover. A fund with the same
leverage but funding locked for months, and margin terms that cannot change overnight, has time; one whose lenders can raise haircuts daily has
the \omterm{def:fm:treasury-and-funding:horizon}{survival horizon} of chapter 14. The two decisions, how much leverage and on what terms, are one decision.
\end{remark}

\section{2006: a natural-gas fund's concentration}

The Senate subcommittee's record describes a hedge fund, Amaranth Advisors, that built very large positions in natural gas contracts: at times 100\,000 contracts in a single month, and, at times in the summer of 2006, about 40\% of all
the outstanding NYMEX contracts for the winter of 2006--2007 and 75\% of those for November 2006. From the end of August to mid-September the natural
gas positions lost more than \$2 billion, and the fund liquidated its entire natural gas portfolio; Senator Coleman's opening statement contrasted it with 1998: the
Federal Reserve did not have to intervene.

\begin{omfigure}
\begin{tikzpicture}
\begin{axis}[width=10.5cm,height=4.8cm,xbar,bar width=10pt,xmin=0,xmax=100,xlabel={\small share of outstanding NYMEX contracts (\%)},
  ytick={0,1,2,3},yticklabels={winter 2006--07,November 2006,January 2007,March 2007},y dir=reverse,ymin=-1.0,ymax=3.6,
  tick label style={font=\footnotesize},grid=major,grid style={gray!20},table/col sep=comma,nodes near coords,nodes near coords style={font=\scriptsize}]
\addplot[fill=omDef!55,draw=omDef] table[y=pos,x=share]{figdata/desk/28-case-studies-i/amaranth.csv};
\draw[gray,dashed] (axis cs:50,-1.0) -- (axis cs:50,3.6) node[pos=0.04,right,font=\scriptsize] {half};
\end{axis}
\end{tikzpicture}
\omcaption{The natural-gas fund's largest shares of the outstanding NYMEX contracts at times in the summer of 2006, from the Senate subcommittee's
record. Data: \texttt{fm\_casebook.AMARANTH}.}\label{fig:fm:case-studies-i:amaranth}
\end{omfigure}

A position that is 75\% of the open interest in a contract has no market to sell into but the other 25\%: the fund is the market. \omterm{def:fm:case-studies-i:days}{Days to liquidate}
makes the point in the units a \omterm{def:fm:the-risk-management-function:cro}{risk committee} uses. With an illustrative average daily volume of 25\,000 contracts in the month's contract, 100\,000
contracts take 20 days at 20\% participation and 40 days at 10\%; a one-day risk measure of such a position describes a risk the firm could not
have exited in a day. Chapter 7's \omterm{def:fm:limits-and-capital-allocation:conc}{concentration limits} exist for this: a limit on a position's share of open interest and on its \omterm{def:fm:case-studies-i:days}{days to liquidate}
binds before a limit on its dollar risk does, because it measures what the position would cost to leave.

\begin{proposition}[Participation and the time to exit]\label{prop:fm:case-studies-i:days}
For a fixed position, \omterm{def:fm:case-studies-i:days}{days to liquidate} is inversely proportional to the participation rate; with a square-root impact law (Book 10), the cost of
exiting at participation $\pi$ grows like $\sqrt\pi$, so halving the cost of the exit requires quartering the participation, which quadruples the days
and the time the position is exposed to the market.
\end{proposition}

\begin{proof}
$D=Q/(\pi V)$ is inversely proportional to $\pi$. Under the square-root law the impact per unit traded at a daily rate $\pi V$ is proportional to
$\sigma\sqrt{\pi V/V}=\sigma\sqrt\pi$, so the total cost of trading $Q$ is proportional to $Q\sigma\sqrt\pi$. Halving it requires $\pi\to\pi/4$, and
then $D\to4D$.
\end{proof}

\section{August 2007: the crowded exit}

In the week of 6 August 2007 a number of quantitative long-short equity hedge funds had unprecedented losses, followed by a significant rebound. On 13
August Goldman Sachs filed a statement that quantitative strategies were under pressure from an increase in overlapping trades, a surge in volatility and an increase in correlations, that its response had been to reduce risk and leverage, and that it and
other investors were putting \$3 billion of equity into one of its funds, Global Equity Opportunities, whose net asset value was about \$3.6 billion.
Khandani and Lo's study of the week, which Book 7 and Book 8 discuss, hypothesised that the losses were initiated by the rapid unwind of one or more
sizable quantitative equity market-neutral portfolios, and that those losses then put pressure on a broader set of equity portfolios.

\begin{omfigure}
\begin{tikzpicture}
\begin{axis}[width=10.5cm,height=5.4cm,xlabel={\small day},ylabel={\small cumulative P\&L (\% of equity)},xmin=0,xmax=20,ymin=-17,ymax=1,
  tick label style={font=\footnotesize},grid=major,grid style={gray!20},table/col sep=comma,legend cell align=left,
  legend style={font=\scriptsize,at={(0.97,0.05)},anchor=south east}]
\addplot[omDef,very thick] table[x=day,y=seller]{figdata/desk/28-case-studies-i/august.csv};
\addplot[omThm,thick,dashed] table[x=day,y=holder]{figdata/desk/28-case-studies-i/august.csv};
\draw[gray,dotted] (axis cs:5,-17) -- (axis cs:5,1) node[pos=0.93,right,font=\scriptsize] {selling ends};
\legend{the fund that sells half its book,a fund that holds}
\end{axis}
\end{tikzpicture}
\omcaption{A synthetic crowded exit: three funds with overlapping books (overlap 0.94) at six times leverage; one sells half its book over five days.
Both lose about 5\% on the first day; the holder loses more by day 5 ($-15.1\%$ against $-13.6\%$) and less by day 20 ($-7.9\%$ against $-9.7\%$),
as the temporary impact decays. Data: \texttt{fm\_cases1.august}.}\label{fig:fm:case-studies-i:august}
\end{omfigure}

The model is Book 7's, and the numbers are synthetic: the point is the shape. The seller's own trades move the prices of the stocks it holds, and
the holder, whose book overlaps it, marks the same moves. Through the selling days the holder loses slightly more, since it still holds its whole
book while the seller's shrinks. When the selling stops, the temporary part of the impact decays and prices partly return: the holder recovers most
of it, while the seller has sold at the bottom of its own move and recovers only on the half it kept. At day 20 the seller is down 9.7\% and the
holder 7.9\%.

What the seller's choice buys is protection against the scenario in which it cannot hold: a lender that calls, a risk limit that forces the cut, an
investor who redeems. A fund with funding that lasts through the episode, and limits set in advance for it, can choose; one without is the seller
whether it chooses or not. The overlap is what turns one firm's decision into everyone's loss: at an overlap of 0.52 the holder loses 4.3\% by day 20,
and with an unrelated book nothing, while the seller's own loss hardly changes.

\section{What the three share}

\begin{method}[Reading a failure for the decisions]\label{met:fm:case-studies-i:read}
\begin{enumerate}
\item Keep to the public record: official reports, court findings, filings. Put every number in a ledger with its source.
\item Reconstruct the firm's capital, positions and funding through the episode from those numbers, stating each derivation.
\item List the decisions taken before the episode that set its terms: leverage, funding term, concentration, limits, who else held the same trades.
\item Rerun the reconstruction with one decision changed, and state what the counterfactual assumes and ignores.
\item Turn the result into a limit or a test the firm can run on itself: a leverage and funding-term pairing, a days-to-liquidate limit, a
crowded-exit stress.
\end{enumerate}
\end{method}

The three cases share a structure. Each firm held positions whose risk looked small in normal markets: spreads between close relatives, calendar
spreads in one commodity, a hedged equity book. Each held them at a size, and with a leverage, that made its exit a market event. And each depended
on others' behaviour: lenders who could raise margins, a market with no one else on the other side, funds with the same trades. The firm-level
decision that recurs is not the trade but the pairing of size, leverage and funding term with the market's capacity to absorb the exit.

\section{Build: the casebook, part 1}\label{bld:fm:case-studies-i:casebook}

\begin{build}
\bfield{Purpose} The public record of the book's case studies as small tables, each figure tied to its ledger row, and the reconstructions and
counterfactuals of chapters 28 and 29.
\bfield{Interface} \texttt{firm.casebook}: \texttt{LTCM}, \texttt{AMARANTH}, \texttt{ltcm\_returns}, \texttt{capital\_path},
\texttt{days\_to\_liquidate}, \texttt{quant\_unwind}; the unwind simulator of \texttt{firm.capacity}.
\bfield{Rules} Only figures from the ledger enter the tables; every derived figure states its derivation; illustrative inputs are named as such.
\bfield{Acceptance tests} \texttt{code/firm/casebook/tests/}: the returns derived from the GAO figures, the capital path at scaled leverage, \omterm{def:fm:case-studies-i:days}{days to
liquidate}, and the seller and holder under a crowded exit.
\bfield{Stretch} A funding-term counterfactual on \texttt{firm.treasury.survival\_horizon}; a margin spiral for the natural-gas fund; the unwind with
several sellers.
\end{build}

\begin{omsources}
\item US General Accounting Office, \emph{Long-Term Capital Management: Regulators Need to Focus Greater Attention on Systemic Risk},
GAO/GGD-00-3, October 1999.
\item US Senate, Permanent Subcommittee on Investigations, \emph{Excessive Speculation in the Natural Gas Market}, hearing, S.~Hrg.~110-235, 2007.
\item The Goldman Sachs Group, Inc., Form 8-K, Exhibit 99.1, 13 August 2007.
\item A.~E. Khandani and A.~W. Lo, ``What happened to the quants in August 2007?'', 2007.
\end{omsources}

\section{Exercises}

\begin{exercise}[$\star$]\label{exo:fm:case-studies-i:1}
From the GAO figures, derive the 1998 fund's capital at the end of July and its return on capital from the end of August to 28 September.
\end{exercise}

\begin{exercise}[$\star$]\label{exo:fm:case-studies-i:2}
Define \omterm{def:fm:case-studies-i:liquidity}{funding liquidity risk} and \omterm{def:fm:case-studies-i:liquidity}{market liquidity risk}, and say which of the two dominates each of the chapter's three cases.
\end{exercise}

\begin{exercise}[$\star$]\label{exo:fm:case-studies-i:3}
How many days does it take to liquidate 100\,000 contracts at 10\% of an average daily volume of 25\,000?
\end{exercise}

\begin{exercise}[$\star\star$]\label{exo:fm:case-studies-i:4}
At 28 to 1, what fall in the value of the balance sheet's assets, net of hedges, wipes out the capital of \$4.67 billion, and what size was the
balance sheet?
\end{exercise}

\begin{exercise}[$\star\star$]\label{exo:fm:case-studies-i:5}
Prove \cref{prop:fm:case-studies-i:days}. What does it imply for a position that is 75\% of the open interest?
\end{exercise}

\begin{exercise}[$\star\star$]\label{exo:fm:case-studies-i:6}
Why does the holder in \cref{fig:fm:case-studies-i:august} lose more than the seller by day 5 and less by day 20?
\end{exercise}

\begin{exercise}[$\star\star\star$]\label{exo:fm:case-studies-i:7}
\emph{Coding.} Rerun \texttt{quant\_unwind} with an overlap of 0.5 and of 0. What are the seller's and the holder's losses at day 20?
\end{exercise}

\begin{exercise}[$\star\star\star$]\label{exo:fm:case-studies-i:8}
\emph{Find the flaw.} ``The position's one-day value at risk is only 5\% of our capital, so it cannot hurt us badly.''
\end{exercise}

\section{Problem: Three Firms, Three Decisions}

\begin{problem}[{Weekend problem --- three firms, three decisions}]\label{pb:fm:case-studies-i:1}
A new \omterm{def:fm:the-risk-management-function:cro}{chief risk officer} is asked by the board which lessons of the best-known quantitative failures apply to the firm, and what to change.

\medskip\noindent\textbf{Part I --- The reading.}
\begin{enumerate}
\item What does a useful reading of a failure keep to, and what does it avoid?
\item Define \omterm{def:fm:case-studies-i:liquidity}{funding liquidity risk} and \omterm{def:fm:case-studies-i:liquidity}{market liquidity risk}.
\item Describe the loop that joins them.
\item Define \omterm{def:fm:case-studies-i:days}{days to liquidate}.
\end{enumerate}

\medskip\noindent\textbf{Part II --- The record.}
\begin{enumerate}[resume]
\item Summarise the public record on the 1998 fund.
\item Why did the Federal Reserve facilitate a private recapitalisation?
\item Summarise the public record on the natural-gas fund.
\item What did Goldman Sachs state on 13 August 2007?
\end{enumerate}

\medskip\noindent\textbf{Part III --- The reconstructions.}
\begin{enumerate}[resume]
\item Derive the 1998 fund's returns on capital by period.
\item Give its capital path, actual and at half the leverage.
\item What does the half-leverage counterfactual assume and ignore?
\item Give the natural-gas fund's \omterm{def:fm:case-studies-i:days}{days to liquidate} at 20\% and 10\% participation.
\item State and prove \cref{prop:fm:case-studies-i:days}.
\item Give the seller's and the holder's losses at days 1, 5 and 20, and explain the crossing.
\end{enumerate}

\medskip\noindent\textbf{Part IV --- The lessons.}
\begin{enumerate}[resume]
\item What do the three cases share?
\item Which limits would you add to the firm's framework?
\item How would you stress a crowded exit on the firm's own book?
\item Why are leverage and funding term one decision?
\item State the \emph{named result}: on the public figures, the 1998 fund's capital at the end of September had it run half the leverage; the
natural-gas fund's \omterm{def:fm:case-studies-i:days}{days to liquidate} at a stated share of volume; and the August 2007 seller's loss against the holder's.
\item In two sentences, write the recommendation.
\end{enumerate}
\end{problem}

\section{Interview questions}

\begin{interviewq}[$\star$ \iqroles{risk}]\label{iq:fm:case-studies-i:1}
What is the difference between \omterm{def:fm:case-studies-i:liquidity}{funding liquidity risk} and \omterm{def:fm:case-studies-i:liquidity}{market liquidity risk}? Give an example of each.
\end{interviewq}

\begin{interviewq}[$\star$ \iqroles{trader, risk}]\label{iq:fm:case-studies-i:2}
A fund runs its balance sheet at 25 times capital. What fall in its assets wipes it out, and why might that overstate its risk?
\end{interviewq}

\begin{interviewq}[$\star\star$ \iqroles{risk}]\label{iq:fm:case-studies-i:3}
One of your traders holds 40\% of the open interest in a futures contract. How do you measure the risk, and what limit would you set?
\end{interviewq}

\begin{interviewq}[$\star\star$ \iqroles{researcher, trader}]\label{iq:fm:case-studies-i:4}
Your statistical-arbitrage book is losing three times its usual daily risk for the third day running and your competitors seem to be too. Do you
sell or hold?
\end{interviewq}

\begin{interviewq}[$\star\star$ \iqroles{risk, trader}]\label{iq:fm:case-studies-i:5}
Why would a fund's creditors recapitalise it rather than let it fail?
\end{interviewq}

\begin{interviewq}[$\star\star\star$ \iqroles{researcher, risk}]\label{iq:fm:case-studies-i:6}
Design a stress test of a crowded exit for a long-short equity book. What inputs does it need that the firm does not observe?
\end{interviewq}
